\chapter{Аналитическая часть}

\section{Постановка задачи}

В соответствии с заданием на курсовую работу необходимо разработать загружаемый модуля ядра ОС Linux для мониторинга изменения приоритетов, времени непрерывной работы, простоя и блокировки процессов. Для решения данной задачи необходимо

\begin{itemize}[label*=---]
	\item проанализировать и выбрать структуры ядра, содержащие необходимую информацию;
	\item проанализировать и выбрать методы передачи информации из модуля ядра в пространство пользователя;
	\item разработать алгоритмы и структуру программного обеспечения;
	\item исследовать приложения различного характера с помощью разработанного программного обеспечения.
\end{itemize}

\section{Планирование в режиме реального времени}

\subsection{Процессы реального времени}

Режим реального времени отражает способность системы обеспечить требуемый уровень сервиса в определённый промежуток времени.
%Работа системы таким образом, что она обрабатывает запросы внешних по отношению к системе процессов (или других систем) за определённый интервал времени.
%Детерминированность --- для выполнения одного сервиса ОС требуется временной интервал с заданной продолжительностью.
Режим реального времени характеризуется временем ответа системы.


К планированию процессов реального времени предъявляются следующие требования: такие процессы должны иметь гарантированное время ответа и никогда не должны вытесняться процессами с более низким приоритетом. Примерами процессов реального времени являются воспроизведение видео и аудио, а также программы, обрабатывающие данные с датчиков. \cite{understanding}

Каждый процесс реального времени имеет приоритет реального времени --- число в диапазоне от 1 до 99. Чем меньше данное значение, тем выше приоритет процесса. Пользователь может изменить приоритет процесса с помощью системных вызовов sched\_setparam() и sched\_setscheduler(). Если несколько процессов реального времени имеют одинаковый приоритет, планировщик выбирает процесс, который стоит первым в очереди готовых процессов. \cite{understanding}

Процесс реального времени вытесняется другим процессом только тогда, когда происходит одно из следующих событий: \cite{understanding}

\begin{itemize}[label*=---]
	\item Процесс вытесняется другим процессом с более высоким приоритетом реального времени.
	\item Процесс выполняет блокирующую операцию и переходит в состояние сна.
	\item Процесс завершается.
	\item Процесс добровольно передаёт управление другому процессу с помощью системного вызова sched\_yield().
	\item Процесс выполняется с политикой планирования SCHED\_RR (Round Robin) и его квант истёк.
\end{itemize}

\subsection{Алгоритмы планирования в режиме реального времени}

Существуют следующие алгоритмы планирования в режиме реального времени.

\begin{itemize}[label*=---]
	\item Rate-Monotonic Scheduling (RMS) --- алгоритм со статическими приоритетами класса планирования. Статические приоритеты назначаются в соответствии с продолжительностью цикла задачи: чем меньше продолжительность цикла, тем выше приоритет процесса. В худшем случае КПД загрузки центрального процессора ограничен величиной
	\begin{equation}
		U=\sum_{i=1}^{n}\frac{C_i}{T_i} \le n(2^{1/n}-1),
	\end{equation}
	где $U$ --- коэффициент использования процессора, $C_i$ --- время выполнения $i$-го процесса, $T_i$ --- период выполнения $i$-го процесса с учётом крайнего срока завершения, $n$ --- число процессов в системе. При этом $\lim\limits_{n \to \infty }U=ln2\approx0.693$. \cite{rms}\cite{habr}
	
	\item Earliest-deadline-first (EDF) Scheduling --- алгоритм с динамическими приоритетами, которые назначаются в соответствии с крайним сроком завершения. Чем раньше крайний срок, тем выше приоритет процесса. В отличие от RMS, планировщик EDF не требует, чтобы процессы были периодическими и постоянно запрашивали одно и то же количество процессорного времени. Единственное требование состоит в том, чтобы каждый процесс объявлял свой крайний срок планировщику, когда он готов к выполнению.\cite{habr}
	\item POSIX real-time-scheduling. Стандарт POSIX.4 определяет три политики планирования. Каждый процесс имеет атрибут планирования, который может быть соответствовать одному из трех вариантов политики. \cite{posix}
	\begin{enumerate}[label*=\arabic*.]
		\item SCHED\_FIFO --- политика упреждающего планирования с постоянным приоритетом, при которой процессы с одинаковым приоритетом обрабатываются по алгоритму FIFO (First in-first out). Данная политика может иметь не менее 32 уровней приоритета.
		\item SCHED\_RR --- политика аналогична SCHED\_FIFO, но использует метод временных срезов (Round-Robin) для планирования процессов с одинаковыми приоритетами. Он также имеет 32 уровня приоритета.
		\item SCHED\_OTHER --- политика зависит от реализации системы.
	\end{enumerate}
\end{itemize}



\section{Анализ структур ядра, предоставляющих информацию о приоритетах, времени выполнения и простоя процессов}

% https://stackoverflow.com/questions/5770770/prio-static-prio-rt-priority-in-linux-kernel
% https://blogs.oracle.com/linux/post/task-priority

\subsection{Структура task\_struct}

В ядре Linux процесс описывается структурой task\_struct. Она определена в файле include/linux/sched.h. \cite{sched} Необходимую информацию о процессе содержат следующие поля:

\begin{enumerate}[label*=\arabic*.]
	\item \textbf{pid} (Process Identifier) --- уникальный идентификатор процесса, по которому можно получить информацию об этом процессе, а также направить ему управляющий сигнал или завершить его.
	\item \textbf{prio} --- значение, которое использует планировщик для принятия решений, связанных с планированием рассматриваемого процесса. Чем ниже значение prio, тем выше приоритет процесса. Константы, задающие диапазон значений приоритетов, представлены в листинге \ref{lst:prio}. \cite{nice}
	
	\begin{lstlisting}[
		caption={Константы, задающие диапазон значений приоритетов процессов, версия ядра 6.1.0},
		label=lst:prio,
		]
#define MAX_NICE	19
#define MIN_NICE	-20
#define NICE_WIDTH	(MAX_NICE - MIN_NICE + 1)

/*
* Priority of a process goes from 0..MAX_PRIO-1, valid RT
* priority is 0..MAX_RT_PRIO-1, and SCHED_NORMAL/SCHED_BATCH
* tasks are in the range MAX_RT_PRIO..MAX_PRIO-1. Priority
* values are inverted: lower p->prio value means higher priority.
*/

#define MAX_RT_PRIO		100

#define MAX_PRIO		(MAX_RT_PRIO + NICE_WIDTH)
#define DEFAULT_PRIO	(MAX_RT_PRIO + NICE_WIDTH / 2)

/*
* Convert user-nice values [ -20 ... 0 ... 19 ]
* to static priority [ MAX_RT_PRIO..MAX_PRIO-1 ],
* and back.
*/
#define NICE_TO_PRIO(nice)	((nice) + DEFAULT_PRIO)
#define PRIO_TO_NICE(prio)	((prio) - DEFAULT_PRIO)
	\end{lstlisting}

	Значение \textbf{nice} (фактор <<любезности>>) --- целое число в диапазоне от -20 до 19 (по умолчанию 0). Чем меньше значение фактора <<любезности>> процесса, тем выше его приоритет.
	
	Диапазон значений prio делится на два интервала:
	\begin{itemize}[leftmargin=1.6\parindent]
		\item от 0 до 99 -- процесс реального времени;
		\item от 100 до 139 -- обычный процесс.
	\end{itemize}

	Реализации функций для вычисления приоритета процесса представлены в листинге \ref{lst:prio_calc}. \cite{prio_calc}
	
	\begin{lstlisting}[
		caption={Реализации функций для вычисления приоритета процесса, версия ядра 6.1.0},
		label=lst:prio_calc,
		]
/* kernel prio, less is more */
static inline int __task_prio(struct task_struct *p)
{
	if (p->sched_class == &stop_sched_class) /* trumps deadline */
		return -2;
	
	if (rt_prio(p->prio)) /* includes deadline */
		return p->prio; /* [-1, 99] */
	
	if (p->sched_class == &idle_sched_class)
		return MAX_RT_PRIO + NICE_WIDTH; /* 140 */
	
	return MAX_RT_PRIO + MAX_NICE; /* 120, squash fair */
}

/*
* Calculate the current priority, i.e. the priority
* taken into account by the scheduler. This value might
* be boosted by RT tasks, or might be boosted by
* interactivity modifiers. Will be RT if the task got
* RT-boosted. If not then it returns p->normal_prio.
*/
static int effective_prio(struct task_struct *p)
{
	p->normal_prio = normal_prio(p);
	/*
	* If we are RT tasks or we were boosted to RT priority,
	* keep the priority unchanged. Otherwise, update priority
	* to the normal priority:
	*/
	if (!rt_prio(p->prio))
		return p->normal_prio;
	return p->prio;
}
	\end{lstlisting}

	\item \textbf{static\_prio} --- значение, которое не изменяется ядром при работе планировщика, однако оно может быть изменено с использованием пользовательского приоритета nice. Реализация функции для изменения значений nice и static\_prio процесса представлены в листинге \ref{lst:static_prio_calc}. \cite{prio_calc}
	
	\begin{lstlisting}[
		caption={Реализация функции для изменения значений nice и static\_prio процесса, версия ядра 6.1.0},
		label=lst:static_prio_calc,
		]
void set_user_nice(struct task_struct *p, long nice)
{
	bool queued, running;
	int old_prio;
	struct rq_flags rf;
	struct rq *rq;
	
	if (task_nice(p) == nice || nice < MIN_NICE || nice > MAX_NICE)
		return;
	/*
	* We have to be careful, if called from sys_setpriority(),
	* the task might be in the middle of scheduling on another CPU.
	*/
	rq = task_rq_lock(p, &rf);
	update_rq_clock(rq);
	
	/*
	* The RT priorities are set via sched_setscheduler(), but we still
	* allow the 'normal' nice value to be set - but as expected
	* it won't have any effect on scheduling until the task is
	* SCHED_DEADLINE, SCHED_FIFO or SCHED_RR:
	*/
	if (task_has_dl_policy(p) || task_has_rt_policy(p)) {
		p->static_prio = NICE_TO_PRIO(nice);
		goto out_unlock;
	}
	queued = task_on_rq_queued(p);
	running = task_current(rq, p);
	if (queued)
		dequeue_task(rq, p, DEQUEUE_SAVE | DEQUEUE_NOCLOCK);
	if (running)
		put_prev_task(rq, p);
	
	p->static_prio = NICE_TO_PRIO(nice);
	set_load_weight(p, true);
	old_prio = p->prio;
	p->prio = effective_prio(p);
	
	if (queued)
		enqueue_task(rq, p, ENQUEUE_RESTORE | ENQUEUE_NOCLOCK);
	if (running)
		set_next_task(rq, p);
	
	/*
	* If the task increased its priority or is running and
	* lowered its priority, then reschedule its CPU:
	*/
	p->sched_class->prio_changed(rq, p, old_prio);
	
out_unlock:
	task_rq_unlock(rq, p, &rf);
}
	\end{lstlisting}

	\item \textbf{normal\_prio} --- значение, которое зависит от статического приоритета и политики планировщика задач. Если политикой планирования процесса является SCHED\_DEADLINE, то его приоритет будет равен -1. Для процессов не реального времени данное значение равняется значению статического приоритета static\_prio, то есть 120 + nice. Для процессов реального времени данное значение равняется значению, вычисленному с использованием максимального значения приоритета процесса реального времени и, непосредственно, его rt\_prio. Важно отметить тот факт, что, чем больше значение rt\_priority процесса, тем выше его приоритет. Функция вычисления нормального приоритета процесса представлена в листинге \ref{lst:normal_prio_calc}. \cite{prio_calc}
	
	\begin{lstlisting}[
		caption={Реализация функции для вычисления нормального приоритета процесса, версия ядра 6.1.0},
		label=lst:normal_prio_calc,
		]
static inline int __normal_prio(int policy, int rt_prio, int nice)
{
	int prio;
	
	if (dl_policy(policy))
		prio = MAX_DL_PRIO - 1;
	else if (rt_policy(policy))
		prio = MAX_RT_PRIO - 1 - rt_prio;
	else
		prio = NICE_TO_PRIO(nice);
	
	return prio;
}

/*
* Calculate the expected normal priority: i.e. priority
* without taking RT-inheritance into account. Might be
* boosted by interactivity modifiers. Changes upon fork,
* setprio syscalls, and whenever the interactivity
* estimator recalculates.
*/
static inline int normal_prio(struct task_struct *p)
{
	return __normal_prio(p->policy, p->rt_priority, 
		PRIO_TO_NICE(p->static_prio));
}
	\end{lstlisting}
	
	\item \textbf{utime} --- это процессорное время, полученное процессом в режиме пользователя.
	\item \textbf{stime} --- это процессорное время, полученное процессом в режиме ядра.
\end{enumerate}



\subsection{Структура sched\_info}

В ядре Linux дескриптор процесса (struct task\_struct) ссылается на структуру sched\_info, предоставляющую информацию о планировании процесса. Структура представлена на листинге \ref{lst:info}. \cite{sched}

\begin{lstlisting}[
	caption={Определение struct sched\_info, версия ядра 6.1.0},
	label=lst:info,
	]
struct sched_info {
	#ifdef CONFIG_SCHED_INFO
	/* Cumulative counters: */
	
	/* # of times we have run on this CPU: */
	unsigned long			pcount;
	
	/* Time spent waiting on a runqueue: */
	unsigned long long		run_delay;
	
	/* Timestamps: */
	
	/* When did we last run on a CPU? */
	unsigned long long		last_arrival;
	
	/* When were we last queued to run? */
	unsigned long long		last_queued;
	
	#endif /* CONFIG_SCHED_INFO */
};
\end{lstlisting}

Поле \textbf{run\_delay} содержит количество времени, которое процесс провёл в очереди на выполнение.



\subsection{Структура sched\_statistics}

В ядре Linux дескриптор процесса (struct task\_struct) ссылается на структуру sched\_statistics, предоставляющую информацию о планировании процесса. Она определена в файле include/linux/sched.h. Ниже перечислены требующиеся в работе поля данной структуры (версия ядра 6.1.0). \cite{sched}

\begin{enumerate}[label*=\arabic*.]
	\item \textbf{wait\_max} --- максимальное время, которое процесс или поток провёл в очереди, ожидая выделение ресурсов.
	\item \textbf{wait\_count} --- общее количество раз, когда процесс или поток попадал в очередь на выделение ресурсов.
	\item \textbf{wait\_sum} --- суммарное время, которое процесс или поток провёл в очереди на выделение ресурсов.
	\item \textbf{iowait\_count} --- количество раз, когда процесс или поток был блокирован в ожидании завершения ввода/вывода.
	\item \textbf{iowait\_sum} --- суммарное время, которое процесс или поток был блокирован в ожидании завершения ввода/вывода.
	\item \textbf{sleep\_max} --- максимальное время, которое процесс или поток провёл в состоянии прерываемого сна.
	\item \textbf{sum\_sleep\_runtime} --- суммарное время, которое процесс или поток провел в состоянии прерываемого сна.
	\item \textbf{block\_max} --- максимальное время, которое процесс или поток провёл в состоянии блокировки.
	\item \textbf{sum\_block\_runtime} --- суммарное время, которое процесс или поток провёл в состоянии блокировки.
	\item \textbf{exec\_max} --- максимальное время, которое процесс или поток провёл в состоянии выполнения.
\end{enumerate}

Представленные выше поля структуры shed\_statistics процесса обновляются во время его выполнения с помощью соответствующих функций, описанных в файле /kernel/sched/stats.h. \cite{stats}



\section{Анализ функций ядра, позволяющих определить принадлежность процесса к процессам реального времени}

\subsection{Функция rt\_prio}

Функция rt\_prio определяет, относится ли процесс к процессам реального времени, с указанием значения приоритета реального времени. Приоритет процесса сравнивается с максимальным значением приоритета для процесса реального времени. Если данное условие истинно, то процесс относится к процессам реального времени и функция возвращает единицу, иначе ноль. Стоит отметить, что макрос unlikely в данном случае применяется для оптимизации выполнения сравнения. Реализация функции rt\_prio приведена в листинге \ref{lst:rt_prio}.~\cite{rt}

\begin{lstlisting}[
	caption={Реализация функции rt\_prio, версия ядра 6.1.0},
	label=lst:rt_prio,
	]
static inline int rt_prio(int prio)
{
	if (unlikely(prio < MAX_RT_PRIO))
		return 1;
	return 0;
}
\end{lstlisting}



\subsection{Функция rt\_task}

Функция rt\_task <<обёртывает>> вызов функции rt\_prio. Если заданный процесс (struct task\_struct) является задачей реального времени (приоритет меньше максимального для задачи реального времени), то возвращается единица, иначе ноль. Реализация функции rt\_task приведена в листинге \ref{lst:rt_task}. \cite{rt}

\begin{lstlisting}[
	caption={Реализация функции rt\_task, версия ядра 6.1.0},
	label=lst:rt_task,
	]
static inline int rt_task(struct task_struct *p)
{
	return rt_prio(p->prio);
}
\end{lstlisting}



\subsection{Функция task\_is\_realtime}

Функция task\_is\_realtime анализирует политику планирования, установленную для заданного процесса (struct task\_struct). Планировщик Linux оперирует следующими политиками планирования:
\begin{itemize}[label*=---]
	\item SCHED\_FIFO --- планирование по алгоритму FIFO (First in-first out).
	\item SCHED\_RR --- планирование по алгоритму Round-Robin.
	\item SCHED\_DEADLINE --- задание максимального срока для выполнения спорадических\footnote{Спорадический (др. греч. sporadikos) --- единичный.} процессов.  Спорадический процесс --- это процесс, состоящий из последовательности работ, каждая из которых активируется не более одного раза за период.  Каждая такая работа имеет срок завершения, и время, необходимое для её выполнения. Данная политика использует алгоритм GEDF (Global Earliest Deadline First).
	\item SCHED\_OTHER --- планирование разделения времени по умолчанию.
	\item SCHED\_BATCH --- планирование пакетных процессов.
	\item SCHED\_IDLE --- планирование низкоприоритетных процессов.
\end{itemize}
SCHED\_FIFO и SCHED\_RR являются политиками реального времени. \cite{policy}

Реализация функции task\_is\_realtime приведена в листинге \ref{lst:task_is_realtime}. \cite{rt}

\begin{lstlisting}[
	caption={Реализация функции task\_is\_realtime, версия ядра 6.1.0},
	label=lst:task_is_realtime,
	]
static inline bool task_is_realtime(struct task_struct *tsk)
{
	int policy = tsk->policy;
	
	if (policy == SCHED_FIFO || policy == SCHED_RR)
		return true;
	if (policy == SCHED_DEADLINE)
		return true;
	return false;
}
\end{lstlisting}





\section{Передача данных из пространства ядра в пространство пользователя}

В ОС Linux для передачи данных из пространства ядра в пространство пользователя используются сокеты netlink и виртуальная файловая система proc, которая предоставляет системные вызовы для реализации взаимодействия между двумя этими пространствами.

Структура proc\_ops определена в файле include/linux/proc\_fs.h и содержит указатели на функции для работы с файлами в виртуальной файловой системе proc. Структура представлена на листинге \ref{lst:proc_ops}. \cite{proc_ops}

\begin{lstlisting}[
	caption={Определение struct proc\_ops, версия ядра 6.1.0},
	label=lst:proc_ops,
	]
struct proc_ops {
	unsigned int proc_flags;
	int	(*proc_open)(struct inode *, struct file *);
	ssize_t	(*proc_read)(struct file *, char __user *, size_t, 
	loff_t *);
	ssize_t (*proc_read_iter)(struct kiocb *, struct iov_iter *);
	ssize_t	(*proc_write)(struct file *, const char __user *, size_t, 
	loff_t *);
	/* mandatory unless nonseekable_open() or equivalent is used */
	loff_t	(*proc_lseek)(struct file *, loff_t, int);
	int	(*proc_release)(struct inode *, struct file *);
	__poll_t (*proc_poll)(struct file *, struct poll_table_struct *);
	long	(*proc_ioctl)(struct file *, unsigned int, unsigned long);
	#ifdef CONFIG_COMPAT
	long	(*proc_compat_ioctl)(struct file *, unsigned int, unsigned 
	long);
	#endif
	int	(*proc_mmap)(struct file *, struct vm_area_struct *);
	unsigned long (*proc_get_unmapped_area)(struct file *, unsigned 
	long, unsigned long, unsigned long, unsigned long);
} __randomize_layout;
\end{lstlisting}

\subsection{Функции copy\_to\_user() и copy\_from\_user()}

Функции copy\_to\_user() и copy\_from\_user(), определённые в файле include/linux/uaccess.h, позволяют передавать данные между пространством ядра и пространством пользователя. Они возвращают количество байт, которые не удалось скопировать. Реализация данных функций приведена в листинге \ref{lst:copy}.~\cite{copy}

\begin{lstlisting}[
	caption={Реализация функции copy\_to\_user, версия ядра 6.1.0},
	label=lst:copy,
	]
static __always_inline unsigned long __must_check
copy_from_user(void *to, const void __user *from, unsigned long n)
{
	if (check_copy_size(to, n, false))
		n = _copy_from_user(to, from, n);
	return n;
}

static __always_inline unsigned long __must_check
copy_to_user(void __user *to, const void *from, unsigned long n)
{
	if (check_copy_size(from, n, true))
		n = _copy_to_user(to, from, n);
	return n;
}
\end{lstlisting}

\subsection{Интерфейс sequence file}

Функции для работы с файлами-последовательностями, определённые в файле include/linux/seq\_file.h, позволяют передавать данные только из пространства ядра в пространство пользователя. Структура seq\_operations определена в файле /include/linux/seq\_file.h и содержит указатели на функции для работы с файлами-последовательностями. Поля структур seq\_file и seq\_operaions представлены в листинге \ref{lst:seq}. \cite{seq}

\begin{lstlisting}[
	caption={Определение struct seq\_file и struct seq\_operaions, версия ядра 6.1.0},
	label=lst:seq,
	]
struct seq_operations;

struct seq_file {
	char *buf;
	size_t size;
	size_t from;
	size_t count;
	size_t pad_until;
	loff_t index;
	loff_t read_pos;
	struct mutex lock;
	const struct seq_operations *op;
	int poll_event;
	const struct file *file;
	void *private;
};

struct seq_operations {
	void * (*start) (struct seq_file *m, loff_t *pos);
	void (*stop) (struct seq_file *m, void *v);
	void * (*next) (struct seq_file *m, void *v, loff_t *pos);
	int (*show) (struct seq_file *m, void *v);
};
\end{lstlisting}

Функции структуры seq\_operations выполняют следующие действия:

\begin{itemize}[label*=---]
	\item start() --- начало передачи данных.
	\item show() --- передача данных.
	\item next() --- увеличение указателя pos.
	\item stop() --- освобождение памяти.
\end{itemize}

Функция stop() может вызываться в одном из следующих контекстов:

\begin{itemize}[label*=---]
	\item после start() --- окончание передачи данных;
	\item после show() --- буфер заполнен и может быть выделен дополнительный объём памяти;
	\item после next() --- данные для передачи закончились и нужно либо завершить работу, либо перейти на вывод другого массива данных.
\end{itemize}



\section*{Выводы из аналитической части}
В результате проведённого анализа для мониторинга были выбраны структуры task\_struct, sched\_info и sched\_statistics. В структуре task\_struct информацию о приоритете процесса содержат поля:
\begin{itemize}[label*=---]
	\item pid;
	\item prio;
	\item static\_prio;
	\item normal\_prio;
	\item utime;
	\item stime.
\end{itemize} 
В структуре sched\_info информацию о количестве времени, которое процесс провёл в очереди на выполнение, содержит поле run\_delay. В структуре sched\_statistics информацию о времени непрерывной работы, простоя и блокировки процесса содержат поля:
\begin{itemize}[label*=---]
	\item wait\_max;
	\item wait\_count;
	\item wait\_sum;
	\item iowait\_count;
	\item iowait\_sum;
	\item sleep\_max;
	\item sum\_sleep\_runtime;
	\item block\_max;
	\item sum\_block\_runtime;
	\item exec\_max.
\end{itemize}

Функция task\_is\_realtime() позволяет определить, является ли процесс процессом реального времени. Для передачи данных из пространства ядра в пространство пользователя и из пространства пользователя в пространство ядра был выбран метод передачи данных с помощью функций виртуальной файловой системы proc.

% В данных структурах определены поля, которые содержат необходимую информацию:
%\begin{itemize}[label*=---]
%	\item task\_struct: поля pid, prio, static\_prio, normal\_prio, utime, stime;
%	\item sched\_info: поле run\_delay;
%	\item sched\_statistics: поля wait\_max, wait\_count, wait\_sum, iowait\_count, iowait\_sum, sleep\_max, sum\_sleep\_runtime, block\_max, sum\_block\_runtime, exec\_max.
%\end{itemize}

% Для определения принадлежности процесса к процессам реального времени будет использоваться функция task\_is\_realtime(), для передачи данных между пространством ядра и пространством пользователя --- функции copy\_to\_user() и copy\_from\_user().
